\section{A local calculus in the lowest Landau level}
\label{sec:locality}

The orbital basis diagonalizes rotations about an axis, but its middle
orbitals extend around an entire latitude.  We first replace that basis by
a redundant family localized in both latitude and longitude.  Its analysis
map embeds the physical fermions into a larger collection of fermionic
modes with a genuine notion of spatial support.  The use of localized
frames to obtain continuum fermionic locality is developed in
\cite{BDN}; we give the finite-sphere construction and estimates needed
here.  We then prove the
propagation and functional-calculus estimates needed to manipulate local
operators.  All constants in this section are independent of the flux and
the particle number.

\subsection{Localized frames and auxiliary fermions}

Fix an axis and write
\begin{equation}
 m=q\sin^2(\theta/2),\qquad u=1+m,\qquad v=1+q-m,
 \qquad w_q(m)=\left(\frac{uv}{q+1}\right)^{1/2}.
 \label{eq:local-scales}
\end{equation}
The area form is $dm\,d\phi$.  Away from the two polar balls of fixed
radius, changing $m$ by $w_q(m)$ or changing $\phi$ by $1/w_q(m)$ moves a
bounded physical distance.  The following elementary comparisons will
also be used at the poles, with $u,v\geq1$:
\begin{equation}
 \frac{u(x)}{u(y)}+\frac{v(x)}{v(y)}
 +\frac{w_q(x)}{w_q(y)}
 \leq C(1+d(x,y))^2.
 \label{eq:local-scale-ratios}
\end{equation}
Indeed $m$ is one half the squared chordal distance from the north pole;
the triangle inequality gives the assertion for $u$, and the south pole
gives it for $v$.  The assertion for $w_q$ follows from their product.

Here is an explicit frame construction.  Choose a nonnegative smooth
function $\eta$, supported in $(-2,2)$ and positive on $[-1,1]$, and a
smooth function $0\leq\chi\leq1$ on $[0,1]$ equal to one on $[0,1/2]$ and zero on
$[2/3,1]$.  For positive integers $k$ use the nonzero functions
\[
 \chi(j/q)\eta(\sqrt{1+j}-k),\qquad
 (1-\chi(j/q))\eta(\sqrt{1+q-j}-k),\qquad 0\leq j\leq q.
\]
Divide each by the square root of the sum of their squares.  The resulting
windows $\xi_a$ satisfy
\begin{equation}
 \sum_a\xi_a(j)^2=1,\qquad
\operatorname{supp}\xi_a\subset I_a,\qquad
 |I_a|\leq Cw_a,\qquad
 |\Delta^l\xi_a(j)|\leq C_lw_a^{-l},
 \label{eq:frame-windows}
\end{equation}
where $I_a$ is an interval of integers, $m_a$ is the corresponding window
center, and $w_a=w_q(m_a)$.  At the first few windows the constants absorb
the bounded widths.  Every index belongs to a bounded number of intervals.
To check the difference bound, extend each expression smoothly in $j$.
On a northern window $k\asymp\sqrt{1+j}\asymp w_q(j)$; each derivative
of $\sqrt{1+j}$ costs at least $w_q(j)^{-1}$, and derivatives of
$\chi(j/q)$ cost $q^{-1}\leq Cw_q(j)^{-1}$.  The denominator is bounded
above and below, and obeys the same estimates.  The southern windows have
the identical verification.

For a window bearing the index $k$, choose the fixed integer
$L_a=16(k+1)$, which is larger than the diameter of the full bump support
before restriction to $[0,q]$.  Thus $L_a\asymp w_a$ and, for a
northern window, its Fourier length is independent of $q$.
For $0\leq\ell<L_a$ set
\begin{equation}
 z_{a\ell}=L_a^{-1/2}\sum_{j=0}^q
       \xi_a(j)e^{2\pi i j\ell/L_a}e_j,
 \qquad x_{a\ell}=(m_a,2\pi\ell/L_a),
 \label{eq:frame-atoms}
\end{equation}
where $e_j$ is the normalized orbital of Section~\ref{sec:prelim}.
Repeated points near the poles are retained as distinct slots.  Their
multiplicities are bounded.  Write $\Lambda_q$ for this set of slots.
Fourier orthogonality and \eqref{eq:frame-windows} give the exact identity
\begin{equation}
 \sum_{a,\ell}|z_{a\ell}\rangle\langle z_{a\ell}|=I_{U_q}.
 \label{eq:tight-frame}
\end{equation}
Thus $\mathsf S_q\psi=(\langle z_x,\psi\rangle)_{x\in\Lambda_q}$ is an isometry
from $U_q$ to $\ell^2(\Lambda_q)$.  The slots have bounded density:
\begin{equation}
 \#\{x\in\Lambda_q:d(x,y)\leq r\}\leq C(1+r)^2.
 \label{eq:slot-volume}
\end{equation}
The window centers have bounded density in radial physical distance,
and their circles carry $O(w_a)$ slots with spacing bounded below except
inside the bounded polar balls.  This proves \eqref{eq:slot-volume}.
Figure~\ref{fig:local-slots} shows the two scales in this construction.

\begin{figure}[t]
\centering
\begin{tikzpicture}[x=1cm,y=1cm,>=Stealth,font=\small]
 \begin{scope}
  \node[anchor=west] at (-.1,2.5) {(a) A window in orbital index};
  \draw[->] (0,.35)--(4.4,.35) node[right] {$j$};
  \draw[blue!65!black,thick]
     (.3,.35)..controls (1,.35) and (1.05,1.65)..(2.15,1.65)
     ..controls (3.25,1.65) and (3.3,.35)..(4,.35);
  \draw[densely dashed,gray] (2.15,.3)--(2.15,1.85);
  \node[above] at (2.15,1.85) {$\xi_a$};
  \draw[<->] (.5,-.05)--(3.8,-.05)
       node[midway,below] {index width $O(w_a)$};
  \node[below] at (2.15,.35) {$m_a$};
  \node[align=center] at (2.15,-.95)
       {$L_a\asymp w_a$ Fourier phases};
 \end{scope}
 \draw[->,thick] (4.95,1.05)--(5.8,1.05);
 \begin{scope}[shift={(8.1,1.05)}]
  \node at (0,1.45) {(b) Slots on a physical latitude};
  \draw[gray!70] (0,0) circle (1.18);
  \draw[gray!55,densely dashed] (-1.08,.48)
      arc[start angle=180,end angle=0,x radius=1.08,y radius=.29];
  \draw[blue!65!black] (-1.08,.48)
      arc[start angle=180,end angle=360,x radius=1.08,y radius=.29];
  \foreach \a in {15,45,...,345}
    {\fill[blue!65!black] ({1.08*cos(\a)},{.48+.29*sin(\a)}) circle (.035);}
  \fill[red!70!black] (.54,.229) circle (.045);
  \draw[red!60!black,densely dashed] (.54,.229) circle (.18);
  \draw[red!35!black,densely dashed] (.54,.229) circle (.31);
  \draw[->,red!60!black] (1.65,-.25)--(.76,.13);
  \node[align=left,anchor=west] at (1.25,-.57)
       {rapidly localized\\physical atom};
  \node[align=center] at (0,-1.95)
       {angular spacing $O(w_a^{-1})$\\physical spacing $O(1)$};
 \end{scope}
\end{tikzpicture}
\caption{The index window has width comparable to $w_a$; its Fourier
phases place a comparable number of slots around the latitude.
Together with unit-order radial spacing, this gives bounded slot
density in physical distance.  The drawing is schematic: dashed
circles indicate rapidly decaying tails, not compact spatial support.}
\label{fig:local-slots}
\end{figure}

We need two quantitative consequences of this construction.  Let $k_y$
be a normalized coherent orbital centered at $y$.

\begin{lemma}[Uniform frame localization]\label{locality:frame}
For every $D$ there is $C_D$ such that
\begin{equation}
 |\langle z_x,k_y\rangle|\leq C_D(1+d(x,y))^{-D}.
 \label{eq:frame-coherent-decay}
\end{equation}
For two frames constructed as above, possibly with different axes, their
atoms satisfy the same bound with $k_y$ replaced by an atom centered at
$y$.  For each fixed integer $J$, \eqref{eq:frame-coherent-decay} also
holds, with constants depending on $J$, for every rotated polar orbital
$e_j(y)$ with $0\leq j\leq J$.
\end{lemma}

\begin{proof}
Choose the axial gauge and azimuth convention in which the coefficient
of $e_j$ in $k_y$ is $e^{ij\phi_y}p_{q,m}(j)^{1/2}$.  Its modulus is
\[
 p_{q,m}(j)^{1/2},\qquad
 p_{q,m}(j)=\binom qj(m/q)^j(1-m/q)^{q-j}.
\]
They obey, after extension by zero outside $[0,q]$,
\begin{equation}
 |\Delta^l p_{q,m}(j)^{1/2}|
 \leq C_{D,l}w_q(m)^{-1/2-l}
       \left(1+\frac{|j-m|}{w_q(m)}\right)^{-D}.
 \label{eq:binomial-symbol}
\end{equation}
Here and below finite differences may shift the argument by a fixed
number of indices, which only changes the constant.  We include a
verification to specify uniformity near the ends.  First put $p=m/q$
and $\sigma^2=qp(1-p)$, so
$w_q(m)^2=1+q\sigma^2/(q+1)\asymp1+\sigma^2$.
Fourier inversion of a binomial law of variance $\sigma^2$ gives
\[
 \sup_j p_{q,m}(j)\leq\frac1{2\pi}\int_{-\pi}^{\pi}
       \bigl(1-4p(1-p)\sin^2(t/2)\bigr)^{q/2}\,dt
       \leq \frac C{1+\sigma}.
\]
For $0<p<1$, tilt the law by $e^{\lambda j}$.  Its new parameter is
$p_\lambda=pe^\lambda/(1-p+pe^\lambda)$, and its variance differs
from $\sigma^2$ by a factor between $e^{-|\lambda|}$ and
$e^{|\lambda|}$.  Therefore the centered log moment generating
function satisfies
$\log\mathbb E e^{\lambda(X-m)}\leq\sigma^2\lambda^2$
for $|\lambda|\leq1/2$.  At $z=j-m$, choose
$\lambda=z/(2(\sigma^2+|z|))$.  The exact tilting identity and the
preceding point-mass bound give
\begin{equation}
 p_{q,m}(j)\leq\frac C{w_q(m)}
       \exp\left(-\frac{c|j-m|^2}{w_q(m)^2+|j-m|}\right).
 \label{eq:binomial-global-tail}
\end{equation}
The cases $p=0,1$ are direct.  This proves every polynomial tail
bound with the required height, uniformly also when $w_q(m)$ is
bounded.

For the derivative gain first assume $w_q(m)^2\geq16(l+1)$, and let
$h=\min(m,q-m)\asymp w_q(m)^2$.  Interpolate the square root of the
binomial probability by gamma functions, writing the resulting
positive function as $F(x)$.  Put $\psi=(\log\Gamma)'$ and
$\psi_1=(\log\Gamma)''$.  On $|x-m|\leq h/2$, the explicit
formulas
\[
 \begin{aligned}
 (\log F)'(x)=\tfrac12\left(\log\frac p{1-p}
          -\psi(x+1)+\psi(q-x+1)\right),\\
 (\log F)''(x)=-\tfrac12\bigl(\psi_1(x+1)+\psi_1(q-x+1)\bigr)
 \end{aligned}
\]
and the positive-real polygamma series give
$|(\log F)'(m)|\leq Cw^{-2}$,
$(\log F)''(x)\leq-cw^{-2}$, and
$|(\log F)^{(k)}(x)|\leq C_kw^{-k}$ for $k\geq2$.
Stirling's inequalities give $F(m)\leq Cw^{-1/2}$.
Differentiating $F=e^{\log F}$ consequently proves
\[
 |F^{(l)}(x)|\leq C_{D,l}w^{-1/2-l}
                    (1+|x-m|/w)^{-D},\qquad |x-m|\leq h/2.
\]
If $|j-m|\leq h/4$, integrate this derivative over the unit cube
that represents $\Delta^lF(j)$; the entire stencil stays in the
displayed interval.  Outside that region, each nonzero stencil
entry has $|j+a-m|\geq cw^2$, $0\leq a\leq l$.
Equation~\eqref{eq:binomial-global-tail} then supplies an
exponential in that distance, which pays both $w^{-l}$ and every
polynomial tail.  This also handles stencils crossing an endpoint.
Finally, if $w^2<16(l+1)$, bound the finite difference by its $l+1$
values and use \eqref{eq:binomial-global-tail}; the factor $w^{-l}$
is now a constant depending only on $l$.  This completes the proof
of \eqref{eq:binomial-symbol}, including the bounded-width regime.

Insert \eqref{eq:binomial-symbol} in \eqref{eq:frame-atoms}.  For windows
whose radial centers are a bounded distance apart, their widths are
comparable; summation by parts $l$ times gives the factor
$(1+w_a|\phi_x-\phi_y|)^{-l}$, with angular difference modulo $2\pi$.
The absolute sum before summation by parts is bounded.  For separated
radial windows, the tail in \eqref{eq:binomial-symbol} first supplies
arbitrary powers of radial distance.  The ratios of the widths cost
only polynomial powers by \eqref{eq:local-scale-ratios}.  Spending
additional tail orders therefore gives the same angular estimate.
Radial distance plus the shorter latitude arc bounds the spherical
distance from above, proving \eqref{eq:frame-coherent-decay}.

The coherent resolution of the identity has density
$(q+1)/(2\pi q)$ with respect to area.  Inserting it between two frame
atoms reduces the frame-to-frame bound to
\[
 \int_{S_q}(1+d(x,z))^{-D-3}(1+d(z,y))^{-D-3}\,dz
 \leq C_D(1+d(x,y))^{-D}.
\]
The inequality follows by dividing the integral into the sets where
$d(x,z)\geq d(x,y)/2$ and where $d(z,y)\geq d(x,y)/2$.
Finally
\[
 |\langle e_j(y),k_z\rangle|^2
 =\binom qj\sin^{2j}(\vartheta/2)\cos^{2(q-j)}(\vartheta/2),
\]
where $\vartheta$ is the angle between $y$ and $z$.  For bounded $j$,
this is bounded by $C_j(1+d(y,z))^{2j}e^{-c d(y,z)^2}$; increasing the
constant covers bounded $q$.  Inserting the coherent resolution once
more proves the last assertion.
\end{proof}

\begin{lemma}[Common frames for nearby fluxes]\label{locality:common-frames}
For each fixed $d_0$, frames at fluxes $q$ and $t$ with
$0\leq q-t\leq d_0$ may be embedded in the same slots so that index
identification $i_{q,t}:e_j^{(t)}\mapsto e_j^{(q)}$, $j\leq t$,
satisfies the exact identity $\mathsf S_t=\mathsf S_q i_{q,t}$.
All preceding estimates remain uniform.  In the northern third the
windows can be chosen from one system independent of the flux; there
one may use the common metric with radial variable $\sqrt{1+m}$.
\end{lemma}

\begin{proof}
Use the windows and Fourier lengths at flux $q$ for both spaces, and
restrict their index sequences to $0\leq j\leq t$ for the smaller
space.  Equation~\eqref{eq:tight-frame} still holds on those indices.
Only a bounded number of end indices have been removed.  The windows
meeting them have bounded width, or have rapidly small tails in the
estimates in which an endpoint extension occurs.  The slot positions
for the two sphere metrics differ by a bounded distance; the constants
may depend on $d_0$.  For all sufficiently large $q$, every window
centered in the northern third has its full support below $q/2$,
because its width is $O(\sqrt q)$.  On these supports the construction
uses only $\eta(\sqrt{1+j}-k)$ and their fixed normalization, and
their Fourier length is the fixed integer $16(k+1)$.  These complete
slot vectors are therefore independent of $q$ in index coordinates.
The buffer between $q/3$ and $q/2$ permits changing the frame outside
this region without changing its northern slots.  Bounded fluxes may
use pairwise padded frames; no long sequence of common slots is needed
there.
\end{proof}

Whenever operators at these two fluxes are subtracted, we first use
$i_{q,t}$ and extend the smaller operator by the identity on the
added orbital modes.  Thus, under the graded decomposition
$\mathcal F_q=\mathcal F_t\,\widehat\otimes\,
\bigwedge\operatorname{span}\{e_{t+1},\ldots,e_q\}$, the padded
operator is $A\widehat\otimes I$.  Extra modes are put in their
vacuum only when comparing states or taking the physical compression.

The fermionic Fock space over $\ell^2(\Lambda_q)$ is a graded tensor
product of two-dimensional slot spaces.  The exterior isometry
$\Gamma(\mathsf S_q)$ identifies $\mathcal F_q$ with the subspace in which
the orthogonal complement of $\mathsf S_qU_q$ is vacant.  A \emph{physical}
operator is an operator in the CAR algebra generated by $\mathsf S_qU_q$,
extended as the identity on the complementary modes.  It preserves
that vacant-complement subspace.  For $X\subset\Lambda_q$, let
$\mathfrak A(X)$ be the CAR algebra of its slots.  Put
$\mathcal N_{\mathrm{slot}}=\sum_{x\in\Lambda_q}c_x^*c_x$.
An operator has number grade $k\in\mathbb Z$ if
$[\mathcal N_{\mathrm{slot}},A]=kA$, and parity
$\epsilon_A\in\{0,1\}$ if
$(-1)^{\mathcal N_{\mathrm{slot}}}A(-1)^{\mathcal N_{\mathrm{slot}}}
=(-1)^{\epsilon_A}A$.  Operators of even parity on disjoint sets
commute; homogeneous operators obey the
graded relation
\[
 [A,B]_{\mathrm{gr}}=AB-(-1)^{\epsilon_A\epsilon_B}BA=0.
\]
All Hamiltonians below have even parity.  Local expansions can be
projected onto any prescribed number grade by averaging the gauge
action of the total slot number $\mathcal N_{\mathrm{slot}}$.
This action preserves every slot algebra and does not increase
approximation errors.  On physical operators its grading agrees with
that of the physical number operator $\mathcal N$.

\subsection{Local norms and summation}

For a center $x$, an operator $A$, and $l>0$, define
\begin{equation}
 \|A\|_{x,l}=\|A\|+
 \sup_{r\geq1}(1+r)^l
       \inf_{A_r\in\mathfrak A(B_r(x))}\|A-A_r\|.
 \label{locality:ball-norm}
\end{equation}
We call $A$ rapidly local at $x$ if this is finite for every $l$,
with bounds uniform in the volume under consideration.  An interaction
is a family $A_x$ integrated against a measure whose mass in every
unit ball is bounded.  A bound of \emph{size} $a(x)$ means
$\|A_x\|_{x,l}\leq C_l a(x)$ for every $l$.  A positive function $a$
is called tempered if, for fixed $C,\kappa$,
\begin{equation}
 a(x)\leq C a(y)(1+d(x,y))^\kappa.
 \label{eq:tempered-weight}
\end{equation}
The constants in estimates involving $a$ may depend on $C,\kappa$.
Powers and products of $u,v,w$ are examples.

Taking near-best approximations at radii $2^k$ and subtracting successive
ones gives a ball decomposition
\begin{equation}
 A=\sum_{k\geq0}A^{(k)},\qquad
 A^{(k)}\in\mathfrak A(B_{2^{k+1}}(x)),\qquad
 \|A^{(k)}\|\leq C_l\|A\|_{x,l}2^{-kl}.
 \label{eq:ball-decomposition}
\end{equation}
Conversely such a decomposition implies the approximation bound, with
any slightly smaller exponent.  We use this elementary equivalence
to pass between operators and interactions on balls.

An observable supported initially on a large ball needs a separate
convention.  A family has \emph{core radius $r$ and amplitude $a$} if
\begin{equation}
 \|A\|\leq Ca,\qquad
 \inf_{A_R\in\mathfrak A(B_R(x))}\|A-A_R\|
       \leq C_Da(1+R)^{-D}\quad(R\geq C_*r),
 \label{eq:local-core-tail}
\end{equation}
where the geometric constant $C_*$ is independent of $D$.
The amplitude may contain a specified power of $r$.  This convention
does not assert that every centered norm \eqref{locality:ball-norm}
is bounded by $Ca$ without a further radius factor.

\begin{lemma}[Shells outside a fixed core]\label{locality:core-shells}
If \eqref{eq:local-core-tail} holds with $r\geq1$, there is a
decomposition $A=\sum_{\nu\geq0}A_\nu$ with
$A_\nu\in\mathfrak A(B_{C_*2^{\nu+1}r}(x))$ and
$\|A_\nu\|\leq C_Da2^{-\nu D}$ for every $D$.
Consequently, for each fixed $B\geq0$,
\begin{equation}
 \sum_{\nu\geq0}(C_*2^{\nu+1}r)^B\|A_\nu\|
       \leq C_Ba r^B.
 \label{eq:local-core-moment}
\end{equation}
\end{lemma}

\begin{proof}
Take approximants at radii $C_*2^\nu r$ and telescope, starting
with the first approximant as the core.  The core norm is $Ca$;
the other differences are bounded by the two adjoining errors in
\eqref{eq:local-core-tail}.  Sum with any $D>B+1$ to obtain
\eqref{eq:local-core-moment}.
\end{proof}

\begin{lemma}[Products, commutators, and weighted sums]
\label{locality:calculus}
The following bounds are uniform over the slot systems above.
\begin{enumerate}
\renewcommand{\labelenumi}{\textup{(\roman{enumi})}}
\item Products of operators rapidly local at the same center, and
adjoints, preserve rapid locality.  If $A,B$ are
centered at $x,y$, respectively, with bounded local norms, then
$\|[A,B]_{\mathrm{gr}}\|\leq C_D(1+d(x,y))^{-D}$ for every $D$.
\item If $A$ is rapidly local at $x$ and a rapidly local interaction
$K_y$ has tempered size $a(y)$, then
$[\int K_y\,dy,A]_{\mathrm{gr}}$ is rapidly local at $x$, of size
$a(x)$ times the local bounds of $A$.
\item Suppose $A$ is even, $\|A\|_{x,l}\leq1$, and $J_y$ has bounded
rapid local norms.  With $D=1+d(x,y)$ and $C_{xy}=[A,J_y]$,
\begin{equation}
 \|C_{xy}\|\leq C D^{-l},\qquad
 \|C_{xy}\|_{x,l}\leq C,\qquad
 \|C_{xy}^*C_{xy}\|_{x,l}\leq C D^{-l}.
 \label{eq:commutator-square-locality}
\end{equation}
\end{enumerate}
For an operator initially supported in $B_r(x)$, the corresponding
core estimates cost a fixed polynomial in $r$.  At distances larger
than a fixed multiple of $r$, any prescribed decay order is available.
\end{lemma}

\begin{proof}
The product estimate follows from
\[
 \|AB-A_rB_r\|
 \leq\|A-A_r\|\|B\|+\|A_r\|\|B-B_r\|.
\]
For the commutator use approximations of radius less than $d(x,y)/3$;
their graded commutator vanishes.  For the weighted assertion expand
both factors as in \eqref{eq:ball-decomposition}.  Terms with disjoint
balls vanish, while intersecting balls of radii $R,S$ require
$d(x,y)\leq R+S$.  The number or measure of their possible centers is
$O((1+R+S)^2)$, and \eqref{eq:tempered-weight} costs at most
$C(1+R+S)^\kappa$.  To obtain an output moment of order $D$, bound
these costs by $C(1+R+S)^{D+\kappa+2}$ and choose input decay order
larger than $D+\kappa+6$.  The two geometric sums then converge.
This argument also permits assigning each output term to the center
of either input; that convention will be used for differences of
interactions.

For \eqref{eq:commutator-square-locality}, the first estimate follows
by the disjoint-ball argument at exponent $l$.  For approximation
radii $r\leq3D$, use zero as an approximant to $C_{xy}$, giving
$r^l\|C_{xy}\|\leq C$.  For $r>3D$, approximate $A$ in $B_{r/3}(x)$
and $J_y$ in $B_{r/3}(y)$; their commutator is supported in $B_r(x)$
and has error $Cr^{-l}$.  Thus the second estimate holds.  Apply the
product estimate with the small norm $\|C_{xy}\|\leq CD^{-l}$ on one
factor to obtain the third estimate.  Finally a ball $B_r(x)$ contains
$O((1+r)^2)$ slots, and tempered ratios within it cost
$O((1+r)^\kappa)$.  The same sums give the stated polynomial core
costs.  For tails beyond $Cr$, the distance from the original ball
is comparable to the tail radius, so increasing a tail order does not
increase the core-volume exponent.
\end{proof}

\begin{corollary}[The physical interactions]\label{locality:basic-interactions}
\label{locality:toeplitz}
The unperturbed Hamiltonian $H_q$ and the Fock lift of any bounded
real scalar symbol $\varphi$ have physical, neutral, Hermitian rapid
interactions of uniform sizes $C_D$ and $C_D\|\varphi\|_\infty$,
respectively.
\end{corollary}

\begin{proof}
Coefficient localization gives operator-norm localization directly.
If $E_R$ is the one-particle projection onto slots in $B_R(x)$, the
CAR identity $\|c(f)\|=\|f\|_2$ gives
\[
 \|c(\mathsf S_q f)-c(E_R\mathsf S_q f)\|
       =\|(I-E_R)\mathsf S_q f\|_2.
\]
For $f=k_x$ or a bounded polar mode at $x$,
Lemma~\ref{locality:frame} and the slot count
\eqref{eq:slot-volume} make the right side at most
$C_D(1+R)^{-D}$ for every $D$; the approximating annihilator is
supported in $B_R(x)$.

The coherent resolution gives
\[
 d\Gamma(T_q(\varphi))=\frac{q+1}{2\pi q}
       \int_{S_q}\varphi(x)c^*(k_x)c(k_x)\,dA(x).
\]
For $H_q$, the highest pair vector is $v_0=e_0\wedge e_1$.
Schur's Lemma on the spin-$(q-1)$ pair space and its dimension
$2q-1$ give, with Haar mass one,
\begin{equation}
 H_q=(2q-1)\int_{SU(2)}B(gv_0)^*B(gv_0)\,dg.
 \label{eq:coherent-pair-resolution}
\end{equation}
The product $B(gv_0)=c(ge_1)c(ge_0)$ is rapidly local at the
rotated pole by Lemmas~\ref{locality:frame} and
\ref{locality:calculus}.  Its phase along the stabilizer disappears
in its square, so the center measure on the sphere has density
$(2q-1)/(2\pi q)$.  This and $(q+1)/(2\pi q)$ are uniformly bounded.
The same two lemmas localize $c^*(k_x)c(k_x)$.  No derivative of
$\varphi$ enters either the representation or its bound.
\end{proof}

\subsection{Propagation and spectral filters}

We next establish that a time filter preserves this local calculus.
The relevant propagation bound is needed only with arbitrary polynomial
decay, which permits a direct truncation argument.  The underlying
commutator iteration is the Lieb--Robinson method
\cite{LR}, in its graded fermionic form \cite{NSY}; its use with time filters is the locality
mechanism behind quasi-adiabatic continuation and spectral flow
\cite{HW,BMNS}.

\begin{lemma}[Propagation of local observables]\label{locality:dynamics}
Let $K$ be self-adjoint and even, with a rapidly local interaction of
uniform size.  Write $\tau_t^K(A)=e^{itK}Ae^{-itK}$.  For each $D$
there are $D'$ and $M_D$, independent of volume, such that
\begin{equation}
 \|\tau_t^K(A)\|_{x,D}
       \leq C_D(1+|t|)^{M_D}\|A\|_{x,D'}.
 \label{eq:polynomial-propagation}
\end{equation}
The assertion holds also for time-dependent generators with the same
uniform bounds.  In particular evolution over a bounded parameter
interval preserves rapid locality.  If $A$ is supported in $B_r(x)$,
the error in approximating $\tau_t^K(A)$ on $B_R(x)$, for
$R\geq C(r+(1+|t|)^4)$, is bounded by
$C_D\|A\|(1+r)^2R^{-D}$.
\end{lemma}

\begin{proof}
Use \eqref{eq:ball-decomposition} to write $K$ as a sum of terms on
balls.  Taking Hermitian and even parts makes each term Hermitian
and even without changing the total or its estimates.  First retain
only terms of diameter at most $\ell$.  For
this truncated interaction define
\[
 M_{ab}=2\sum_{X\ni a,b}\|K_X\|.
\]
Uniform summability, including the number of sites in $X$, gives
$\sup_a\sum_b M_{ab}e^{d(a,b)/\ell}\leq C$; indeed each retained
distance is at most $\ell$, and the ball-volume costs are polynomial.
The iterated commutator inequality, with internal evolution removed
unitarily at each step, therefore gives
\begin{equation}
 \|[\tau_t^{K_{\leq\ell}}(A),B]_{\mathrm{gr}}\|
 \leq C\|A\|\|B\|\,|X|\,|Y|
        e^{C|t|-d(X,Y)/\ell}
 \label{eq:truncated-LR}
\end{equation}
for $A\in\mathfrak A(X)$ and $B\in\mathfrak A(Y)$.
For completeness, the $n$th iterated integral contributes $|t|^n/n!$
times a chain of $n$ matrix factors $M$.  Multiplication by
$e^{d(X,Y)/\ell}$ is bounded by the product of the exponential
weights along the chain.  Summing the intermediate sites gives
$C^n$, and summing $n$ gives \eqref{eq:truncated-LR}.  Evenness of
the generator makes the same calculation valid for odd $A,B$ with
the graded commutator.

Restore the removed terms using Duhamel's formula.  For each removed
ball term use the smaller of the trivial commutator bound and
\eqref{eq:truncated-LR}.  The centers within distance
$C\ell(1+|t|)$ of $X$, together with the radius of that term, occupy
at most a polynomial volume.  Outside that region the exponential
bound is summable.  Consequently, for every $s$, the restoration
error is at most
\begin{equation}
 C_s\|A\||X|(1+|t|)
       (1+r+\ell(1+|t|))^4\ell^{-s},
 \label{eq:range-restoration}
\end{equation}
when $X\subset B_r(x)$.  To see the power allowance explicitly, a
removed term of radius $a$ contributes its norm times at most
$C(1+r+\ell(1+|t|)+a)^2$ possible centers and at most
$C(1+a)^2$ sites.  Its arbitrarily high radius moment makes the sum
over $a>\ell/2$ bounded by the right side of
\eqref{eq:range-restoration}, after increasing the input moment by
$s+6$.

Now let $R\geq C(r+(1+|t|)^4)$ and choose $\ell=R^{1/3}$.
The truncated evolution differs from the evolution generated by its
terms contained in $B_R(x)$ by
$C\|A\||X|\operatorname{poly}(R)e^{C|t|-cR^{2/3}}$.
This is again Duhamel's formula and \eqref{eq:truncated-LR}, summed
over the boundary terms.  Both restoration errors are bounded by
$C_D\|A\||X|R^{-D}$, choosing $s$ sufficiently large in
\eqref{eq:range-restoration}.  This supplies an approximation in
$\mathfrak A(B_R(x))$.  At smaller radii use the norm bound
$\|\tau_t(A)\|=\|A\|$.  Since $|X|\leq C(1+r)^2$, this proves the
assertion for supported $A$ and then, by
\eqref{eq:ball-decomposition}, \eqref{eq:polynomial-propagation}.
The proof uses only uniform bounds on the interaction at each time,
so also proves the time-dependent assertion.
\end{proof}

\begin{lemma}[Time filters]\label{locality:filters}
Suppose $F\in L^1(\mathbb R)$ and
$\int(1+|t|)^j|F(t)|\,dt<\infty$ for every $j$.  Under the hypotheses
of Lemma~\ref{locality:dynamics},
\begin{equation}
 \mathcal T_{F,K}(A)=\int_{\mathbb R}F(t)\tau_t^K(A)\,dt
 \label{eq:time-filter}
\end{equation}
is rapidly local whenever $A$ is, with uniform bounds.  The same is
true for a finite measure consisting of such a function and a mass at
zero.  Number grade and physicality are preserved when $K$ is neutral
and physical.  Suitable symmetry of $F$ preserves adjoints.
\end{lemma}

\begin{proof}
Integrate \eqref{eq:polynomial-propagation}.  The approximation norm
is a seminorm up to its operator-norm summand, so the integral of local
approximants gives the required estimate.  The remaining assertions
follow directly by commuting the evolution with gauge transformations,
the physical algebra, and adjoints.
\end{proof}

The frequency functions used below are smooth, vanish on an interval
about zero, and have tails proportional to $1/\omega$ on one or both
half-lines.  Their inverse Fourier transforms satisfy the hypotheses
of Lemma~\ref{locality:filters}.  Here is a convenient verification.
Every positive-order derivative of such a multiplier is integrable,
so integration by parts gives arbitrary powers of decay for $|t|\geq1$.
For $0<|t|<1$, split the frequency integral at $|\omega|=|t|^{-1}$.
The inner portion is $O(1+|\log|t||)$, and one integration by parts
on the outer portion is $O(1)$.  Thus the possible singularity at
zero is integrable.  A smooth compactly supported multiplier has a
Schwartz kernel.  These facts cover both spectral transport and
the filters that retain only small energy differences.

\begin{lemma}[Weighted differences]\label{locality:difference}
Let $K,K'$ satisfy Lemma~\ref{locality:dynamics} uniformly, and assume
$K-K'$ has a rapidly local interaction of tempered size $a(z)$.
For a rapidly local $A$ centered at $x$,
\begin{equation}
 \|\tau_t^K(A)-\tau_t^{K'}(A)\|_{x,D}
 \leq C_Da(x)(1+|t|)^{M_D}\|A\|_{x,D'}.
 \label{eq:weighted-evolution-difference}
\end{equation}
The same conclusion holds after an all-moments time filter, and for
differences of bounded-parameter evolutions.  If the input itself
changes by local size $b(x)$, its separate contribution is $Cb(x)$.
All estimates retain the product of a local input size and $a(x)$.
If $A\in\mathfrak A(B_r(x))$, put $s=\kappa+2$, with $\kappa$ from
\eqref{eq:tempered-weight}.  The difference has operator norm at most
\begin{equation}
 C a(x)\|A\|(1+r)^s|t|(1+|t|)^{4s}.
 \label{eq:difference-core}
\end{equation}
For $R\geq C_*(1+r)$, with a fixed geometric constant $C_*$, it has
an approximation supported in $B_{r+R}(x)$
with error at most
\begin{equation}
 C_Da(x)\|A\|(1+r)^s
        (1+|t|)^{4D+4s+1}(1+R)^{-D}.
 \label{eq:difference-tail}
\end{equation}
Thus the power of the initial radius is independent of the tail order.
If the differing interaction vanishes throughout a ball of radius $d$
about $x$, the bounds for a fixed rapidly local input gain arbitrary
inverse powers of $1+d$.
\end{lemma}

\begin{proof}
Use the exact formula
\[
 \tau_t^K(A)-\tau_t^{K'}(A)
 =i\int_0^t\tau_s^K
       ([K-K',\tau_{t-s}^{K'}(A)])\,ds.
\]
Apply Lemma~\ref{locality:dynamics} to the inside observable,
Lemma~\ref{locality:calculus}(ii) to the commutator, and
Lemma~\ref{locality:dynamics} once more to the outside evolution.
The tempered inequality \eqref{eq:tempered-weight} is precisely what
allows the insertion at $z$ to be charged to $a(x)$.  All time costs
are polynomial and therefore integrable against the specified
filters.  The time-dependent version follows from its identical
Duhamel formula.  Linearity handles a changed input.

We give the additional radius bookkeeping for
\eqref{eq:difference-core}--\eqref{eq:difference-tail}.
For a ball decomposition about $x$, use the weighted sum
$\sum_R(1+R)^p\|A_R\|$.  Lemma~\ref{locality:dynamics} bounds its
evolved value by $C_p(1+|t|)^{4p}$ times its initial value.  Indeed,
for a radius-$r$ input, begin a dyadic decomposition at
$Q=C(1+r+(1+|t|)^4)$.  The core contributes $C\|A\|Q^p$;
the subsequent shells contribute at most
$C_M\|A\|(1+r)^2\sum_{R\geq Q}R^{p-M}$.
Choosing $M>p+2$ bounds this by $C_p\|A\|Q^p$.
The weighted insertion estimate in the proof of
Lemma~\ref{locality:calculus} costs $s=\kappa+2$ powers of the input
radius.  Duhamel at $p=0$ proves \eqref{eq:difference-core}.

For the tail first take $|t|\leq cR^{1/4}$ in the Duhamel integrand,
with $c>0$ a sufficiently small fixed constant; bounded $R$ are
covered by \eqref{eq:difference-core}.  Fix also a small $\varepsilon>0$.
Choose $\varepsilon,c$ and then $C_*$ so that the propagation lemma
applies both from radius $r$ to $r+\varepsilon R$ and from radius
$r+3\varepsilon R$ to $r+R$.  This requires only fixed constants:
if $C$ is the constant in that lemma, take
$\varepsilon<(100C)^{-1}$, $c^4<\varepsilon/(10C)$, and
$C_*>100C/\varepsilon$.
Cut the inner evolved observable at radius $r+\varepsilon R$.  Its remaining
shells have weighted sum of order $s$ at most
$C_M\|A\|(1+r)^2R^{s-M}$, by the short-time bound in
Lemma~\ref{locality:dynamics}; their commutator with the insertion
has the same bound times $Ca(x)$.  In the retained part, discard
insertion terms of radius greater than $\varepsilon R$.  Their arbitrary
radius moments give an error $C_Ma(x)\|A\|R^{-M}$.
The remaining commutator is supported in $B_{r+3\varepsilon R}(x)$ and has
norm at most $Ca(x)\|A\|R^s$.  Evolving it in the outside factor
of Duhamel and using the short-time propagation estimate gives
an approximation on $B_{r+R}(x)$ with error
$C_Ma(x)\|A\|R^{s+2-M}$.
Choose $M$ in each of these three bounds larger than the requested
order plus $s+4$, and integrate over an interval of length at most
$R^{1/4}$.  This gives the stated tail bound for short times.
For $|t|>cR^{1/4}$, use \eqref{eq:difference-core} and
$(1+|t|)^{4D}(1+R)^{-D}\geq c_D$.

Finally suppose an insertion ball of radius $a$ can occur only if
its center $z$ satisfies $d(x,z)+a\geq d$.  If that ball meets an
input shell of radius $R$, then $d\leq R+2a$.  Inserting
$(1+d)^{-M}(1+R+2a)^M$ into the shell summation proves the last
assertion, after increasing the input moments.  This includes changes
supported outside a common comparison region.
\end{proof}

To specify physical compression, let $S:U_q\to K_S$ and
$T:U_q\to K_T$ be two slot analysis isometries.  Write $J_S=\Gamma(S)$
for the vacuum-complement embedding and $\iota_T$ for the embedding of
the physical CAR algebra into the output slots, acting as the identity
on complementary modes.  The map that we estimate is
\begin{equation}
 \mathcal C_{T\leftarrow S}(A)
   =\iota_T(J_S^*AJ_S).
 \label{eq:physical-compression-map}
\end{equation}
When $T=S$ this is compression followed by identity extension; it is
a unital contractive completely positive map.  On physical input
operators it is simply change of frame.

\begin{lemma}[Compression, changes of frame, and smooth functions]
\label{locality:functional-calculus}
The map \eqref{eq:physical-compression-map}, for the frames of
Lemma~\ref{locality:frame}, preserves rapid locality.  An
operator supported in $B_r(x)$ remains localized about that ball,
with polynomial core costs and arbitrary inverse-power tails outside
$B_{Cr}(x)$.  If $A=A^*$ is bounded and rapidly local, and $f$ is
smooth on a neighborhood of its spectrum, then $f(A)$ is rapidly
local.  The constants depend on finitely many local bounds and
derivatives for each requested decay order.
\end{lemma}

\begin{proof}
Put two slot spaces side by side.  If $S,T$ are their analysis
isometries, the partial isometry $TS^*$ has rapidly decaying matrix
entries by Lemma~\ref{locality:frame}.  On their direct sum, the
anti-self-adjoint one-particle operator
\[
 \begin{pmatrix}0&-ST^*\\TS^*&0\end{pmatrix}
\]
rotates the two physical subspaces into each other through angle
$\pi/2$ and acts as zero on their complements.  Its second
quantization is a bounded-strength, rapidly local quadratic
interaction: absolute row sums and every distance moment are bounded
by \eqref{eq:slot-volume}.  Lemma~\ref{locality:dynamics} applies
over this fixed angle.  Conjugate $A\otimes I$ by this rotation and
take the vacuum expectation in the entire first slot copy, retaining
the entire second copy.  On a CAR monomial, the physical components
have moved from $S U_q$ to $T U_q$, while the complementary components
in the first copy are evaluated in their vacuum.  The resulting map
is exactly \eqref{eq:physical-compression-map}, including for odd
monomials; linearity gives the identity for every operator.  Vacuum
expectation is contractive and preserves support in the doubled slot
geometry.  This proves the assertion, and $S=T$ gives re-embedded
physical compression.  In particular no projection onto the output
complementary vacuum is left in the resulting operator.

For the functional-calculus assertion extend $f$ smoothly to a compactly
supported function on $\mathbb R$.  If $A_r$ approximates $A$ locally,
its Hermitian part does so with the same error.  The Fourier formula
and Duhamel's identity give
\[
 \|f(A)-f(A_r)\|
 \leq\left(\int_{\mathbb R}|t|\,|\widehat f(t)|\,dt\right)
       \|A-A_r\|.
\]
Since $f(A_r)$ belongs to the same local algebra, the approximation
bounds follow.  A constant value of $f$ contributes only a scalar
identity, which is local in every ball.
\end{proof}

\subsection{Pair rows and diagonal symbols}

The coherent orbit representations already prove locality of the
unperturbed Hamiltonian and bounded scalar perturbations.  Pinning zeros
will deform the pair coefficients and remove full rotation invariance.
To control those deformations and compare neighboring fluxes, we now
prove a local calculus directly for angular-momentum rows and diagonal
orbital symbols.

Write $Q_b=B_{b-1}$ for the normalized pair annihilator with index
sum $b$ in \eqref{eq:pair-rows}, $1\leq b\leq2q-1$.  Define its
ordered coefficients by
\begin{equation}
 a_b(i)=\frac{(b-2i)\sqrt{\binom qi\binom q{b-i}}}{\sqrt{S_b}},
 \qquad
 S_b=\binom{2q}{b}\frac{b(2q-b)}{2q-1}.
 \label{eq:pair-symbol-normalization}
\end{equation}
In this formula the sum for $S_b$ is over ordered $i$.  With a fixed
choice of annihilator sign,
$Q_b=2^{-1/2}\sum_{i+j=b}a_b(i)c_jc_i
=\sqrt2\sum_{i<j,\,i+j=b}a_b(i)c_jc_i$.
Changing this common sign leaves every square unchanged.

\begin{lemma}[Pair profiles and local row squares]
\label{locality:row-symbols}
With $m=b/2$ and $w=w_q(m)$, the zero-extended coefficients satisfy
\begin{equation}
 |\Delta_i^r\Delta_b^s a_b(i)|
 \leq C_{D,r,s}w^{-1/2-r-s}
       \left(1+\frac{|i-m|}{w}\right)^{-D}.
 \label{eq:pair-row-symbols}
\end{equation}
There are pair annihilators $L_{\nu,\phi}$, rapidly local with uniform
bounds about centers $x_{\nu,\phi}$, and a measure $d\mu$ of bounded
mass per unit ball, such that
\begin{equation}
 H_q=\sum_\nu\int_0^{2\pi}
                L_{\nu,\phi}^*L_{\nu,\phi}\,d\mu_\nu(\phi).
 \label{eq:pair-local-resolution}
\end{equation}
The construction is linear in the row coefficients.  In particular
symbol bounds \eqref{eq:pair-row-symbols} with an additional tempered
factor give local bounds with that factor; the same holds for
differences of row families in common frames.
\end{lemma}

\begin{proof}
Vandermonde's identity identifies
$\binom qi\binom q{b-i}/\binom{2q}{b}$ as a hypergeometric law of
mean $b/2$ and variance $b(2q-b)/(4(2q-1))$.
Multiplying its second centered moment by four proves
\eqref{eq:pair-symbol-normalization}.  The moment formula itself
also follows directly by applying
$i\binom qi=q\binom{q-1}{i-1}$ twice and Vandermonde's identity.

Here are symbol estimates that include the dependence on the moving
row center.  Set
\[
 \ell_q(x)=\log\Gamma(q+1)-\log\Gamma(x+1)
                         -\log\Gamma(q-x+1),\quad t=i-m,
\]
and
\[
 E(m,t)=\tfrac12\{\ell_q(m+t)+\ell_q(m-t)\}-\ell_q(m),
 \qquad T(m)=S_{2m}e^{-2\ell_q(m)},
\]
where the gamma formula defines the smooth extension.  Then
$a_b(i)=-2t e^{E(m,t)}/\sqrt{T(m)}$.  The gamma duplication formula
gives the exact expression
\begin{equation}
 T(m)=\frac{4m(q-m)}{2q-1}\sqrt\pi\,
       \frac{R(m)R(q-m)}{R(q)},
 \qquad R(x)=\frac{\Gamma(x+1)}{\Gamma(x+1/2)}.
 \label{eq:pair-moving-normalization}
\end{equation}
For $m,q-m\geq1/2$, Stirling's inequalities yield
$T(m)\asymp w_q(m)^3$.  Its logarithmic derivatives of order $r\geq1$
are bounded by $C_r(u^{-r}+v^{-r})$.  These derivative assertions
can be obtained directly by differentiating Euler's gamma product:
for $k\geq1$,
\[
 \frac{d^{k+1}}{dx^{k+1}}\log\Gamma(x)
       =(-1)^{k+1}k!\sum_{n=0}^{\infty}(x+n)^{-k-1}.
\]
In the region $|t|\leq c\min(u,v)$, the second derivative
$\ell_q''$ is negative and comparable in magnitude to
$u^{-1}+v^{-1}\asymp w^{-2}$.  Thus
$E(m,t)\leq-ct^2/w^2$.  Differentiating $E$ in $i$ and $b$ gives
polynomials in $t/w$, with one factor $w^{-1}$ per derivative;
\eqref{eq:pair-moving-normalization} has the same property.
The discrete product rule, or integration of these derivative bounds
over unit cubes, proves \eqref{eq:pair-row-symbols} in this region.
There is also a global Gaussian bound, including the ends: convexity
of the trigamma function $\psi_1=(\log\Gamma)''$ and
$\psi_1(x)\geq1/x$ give
\[
 E_{tt}(m,t)=-\tfrac12\bigl\{
 \psi_1(m+t+1)+\psi_1(m-t+1)
 +\psi_1(q-m+t+1)+\psi_1(q-m-t+1)\bigr\}
 \leq-\frac1u-\frac1v\leq-\frac c{w^2}.
\]
Since $E(m,0)=E_t(m,0)=0$, we have $E(m,t)\leq-ct^2/w^2$
on the entire allowed interval.  Outside the central region
$|t|/w\geq cw$ when $w$ is large, so increasing the tail order pays
all derivative and endpoint-extension factors.  The rows with bounded
$w$ require only bounded finite differences.  This proves the full
uniform estimate without differentiating a lattice normalization sum.

Choose smooth windows $\zeta_\nu(b)$ with
$\sum_\nu\zeta_\nu(b)^2=1$, supported on intervals of length
$O(W_\nu)$ around $b=2M_\nu$, where $W_\nu=w_q(M_\nu)$.
The construction \eqref{eq:frame-windows}, with $m=b/2$, supplies
them.  Define
\begin{equation}
 L_{\nu,\phi}=W_\nu^{-1/2}
        \sum_b\zeta_\nu(b)e^{-ib\phi}Q_b,
 \qquad d\mu_\nu(\phi)=W_\nu\frac{d\phi}{2\pi}.
 \label{eq:fourier-pair-localization}
\end{equation}
Fourier orthogonality proves \eqref{eq:pair-local-resolution} exactly.
The radial centers have bounded density in physical distance, while
their circle lengths are comparable to $W_\nu$ outside the polar
balls.  Hence $\mu$ has bounded mass on unit balls.

After removing its linear phase $e^{-i(i+j)\phi}$, the two-index
coefficient of $L_{\nu,\phi}$ has height $O(W_\nu^{-1})$,
and every mixed finite difference of orders $r,s$ costs
$W_\nu^{-r-s}$.  Its tails are arbitrary powers of
$1+|i-M_\nu|/W_\nu+|j-M_\nu|/W_\nu$.
To express it in slot fermions, multiply in each index by the window
coefficients from \eqref{eq:frame-atoms} and sum.  On comparable
windows, repeated summation by parts in both indices gives, for arbitrary $D$, coefficients
bounded by
\[
 C_D(1+d(x,x_{\nu,\phi}))^{-D}
       (1+d(y,x_{\nu,\phi}))^{-D}
\]
for the slot monomial $c_xc_y$.  The height is bounded because the
two summation lengths are $O(W_\nu)$, the coefficient height is
$O(W_\nu^{-1})$, and the two Fourier normalizations contribute
$O(W_\nu^{-1})$.  Noncomparable windows have radial separation;
their scale ratios cost polynomial powers by
\eqref{eq:local-scale-ratios}, paid by increasing the radial tail
orders.  This proves the displayed bound for all slots.  Its absolute
sum, including every distance moment, is bounded by
\eqref{eq:slot-volume}.  Since $\|c_xc_y\|\leq1$, truncating that
sum proves rapid locality in operator norm.  Products give locality
of $L_{\nu,\phi}^*L_{\nu,\phi}$.  Each operation used was linear in
the row coefficients before the final square; keeping any extra
tempered factor proves the last assertion.
\end{proof}

For a smooth real function $a$ on $[0,q]$, fix a constant $C_0$
larger than the window-support constants and define
\begin{equation}
 \mathfrak s_K(a;m)=\max_{0\leq k\leq K}
   \sup_{\substack{j\in[0,q]\\|j-m|\leq C_0w_q(m)}}
       w_q(m)^k|a^{(k)}(j)|.
 \label{eq:diagonal-symbol-seminorm}
\end{equation}
The order-zero term is part of this seminorm.  Scaled derivatives,
rather than merely bounded unscaled derivatives, are necessary:
orbital oscillations on the scale of one index translate an operator
around a macroscopic latitude.

\begin{lemma}[Diagonal symbols]\label{locality:one-body-symbol}
Suppose $A(m)>0$ is tempered with exponent $\kappa$ and
$\mathfrak s_K(a;m)\leq C_KA(m)$ for every integer $K\geq0$.
Then $D_q(a)=\sum_j a(j)c_j^*c_j$ has a physical, neutral, Hermitian
interaction of size $A(m)$.  For each output approximation order $D$,
only the seminorms through
$K_D=\lceil4D+4\kappa+40\rceil$ are needed.  The construction and
the estimate are linear in the symbol, so the same assertion holds
for differences of symbols.
\end{lemma}

\begin{proof}
Insert the analysis map on both sides of the diagonal one-particle
matrix.  The slot matrix elements are
\[
 (L_{a_1}L_{a_2})^{-1/2}\sum_j
   a(j)\xi_{a_1}(j)\xi_{a_2}(j)
                  e^{ij(\phi_2-\phi_1)}.
\]
The two radial windows must overlap, so their widths are comparable
and their radial physical separation is bounded.  Taking $K$
finite differences of the summand and summing by parts gives the
bound $C_K A(m_1)(1+d(x_1,x_2))^{-K}$.  The normalization cancels
the length of the summation interval.  At bounded-width end windows
the same statement follows from finite sums; all involved slots are
inside a bounded polar ball.

Assign the quadratic monomial to $x_1$, and pair it with its adjoint
to obtain Hermitian terms.  Vacuum-compress each assigned term using
\eqref{eq:physical-compression-map}.  Compression fixes the physical
total $D_q(a)$, so the resulting decomposition is still exact and
now each term is physical.  A monomial joining points at distance
$R$ is supported in $B_{1+R}(x_1)$; the compression lemma bounds its
approximation norm of order $D$ by $C_D(1+R)^{D+2}$.
The slot count \eqref{eq:slot-volume}, the tempered ratios, and
$K>D+\kappa+6$ therefore make the assigned absolute sum converge
with the required $D$th moment.  The stated value of $K_D$ more than
covers these losses and the fixed finite differences of the windows.
The grade and reality are preserved by compression.  Every step is
linear before pairing adjoints, which also proves the difference
assertion.
\end{proof}

We record explicitly how the symbol criterion applies after subtracting
an affine function.  Let $\rho(j,m)$ denote meridional distance between
the latitudes of indices $j,m$, and fix a center $x$ at index $m$.
Suppose there are $\beta>0$, $\kappa\geq0$, and constants $C_k$ such
that, for every $k\geq2$ and $j\in[0,q]$,
\begin{equation}
 w_q(j)^k|a^{(k)}(j)|
       \leq C_k\beta(1+\rho(j,m))^\kappa.
 \label{eq:affine-symbol-hypothesis}
\end{equation}
Define $\widetilde a_x(j)=a(j)-a(m)-a'(m)(j-m)$.
Taylor's formula and \eqref{eq:local-scale-ratios} imply
\begin{equation}
 \mathfrak s_K(\widetilde a_x;m_y)
       \leq C_K\beta(1+d(x,y))^{\kappa+8}.
 \label{eq:affine-symbol-size}
\end{equation}
Here is the bound for the two orders not already present in
\eqref{eq:affine-symbol-hypothesis}.  The index excursion between
$m$ and $j$ is at most a constant times
$w_q(j)(1+\rho(j,m))^2$, and the ratio of $w_q(j)$ to the minimum
of $w_q$ on that meridian segment is polynomially bounded by
$(1+\rho(j,m))^2$.  Integrating $a''$ once or twice therefore bounds
$w_q(j)|\widetilde a_x'(j)|$ and $|\widetilde a_x(j)|$ by the
right side of \eqref{eq:affine-symbol-size}.  On each window
the widths are comparable and meridional distance changes by a
bounded amount, proving the displayed seminorm bound.

Lemma~\ref{locality:one-body-symbol} and the weighted commutator
estimate now show that $[D_q(\widetilde a_x),O_x]$ has rapid local
size $C\beta$ whenever $O_x$ has bounded rapid local norms at $x$.
In applications one can take $\beta$ comparable to
$w_q(m)^2|a''(m)|$ and verify
\eqref{eq:affine-symbol-hypothesis} directly.  For an annular sum
commuting with $\mathcal N$ and $\sum_jj c_j^*c_j$, replacing $a$
by $\widetilde a_x$ leaves the commutator unchanged: the removed
operator is a linear combination of those two conserved quantities.

Commutators, evolution, and spectral functions of physical operators
remain physical, and neutral evolution preserves number grade.  The
strictly supported approximants used to measure locality need not
preserve the vacant-complement subspace; physical vacuum compression
is available when a physical representative is required.
The frame may therefore be fixed while proving an operator estimate,
changed when rotating an observable, or chosen compatibly at nearby
fluxes.  Lemmas~\ref{locality:filters} and~\ref{locality:difference}
ensure that the same local estimates remain available after the
spectral transports used next.
